\subsection{Uniqueness of integral representations}

Let E be a Hausdorff weak locally convex space (TVS, II, §6, No.~2), C a proper pointed convex cone in E. One knows that C is the set of elements \(\geqslant0\) of E for an order relation compatible with the vector space structure of E. When C is said to be lattice-ordered, it is of course the order induced on C by that of E that is understood.

Lemma 5. --- Assume that C is weakly complete. Let A be the set of restrictions to C of the continuous linear forms on E. Let \((f_{\lambda})_{\lambda \in \Lambda}\) be a finite family of elements of A, and \(f = \sup(f_{\lambda})\) . For every \(x \in C\) , define

\[
\overline {{{f}}} (x) = \sup \left(f (x _ {1}) + f (x _ {2}) + \dots + f (x _ {n})\right),
\]

the supremum being taken over the set \(S_{x}\) of sequences \((x_{1}, x_{2}, \ldots, x_{n})\) of elements of C such that \(x_{1} + x_{2} + \cdots + x_{n} = x\) . Let \(\operatorname{Card}\Lambda = p\) . Then there exists \((y_{1}, \ldots, y_{p}) \in \operatorname{S}_{x}\) such that \(\overline{f}(x) = f(y_{1}) + \cdots + f(y_{p})\) .

Denote by \(f_1, f_2, \ldots, f_p\) the elements of the family \((f_\lambda)\) . For \(k = 1, 2, \ldots, p\) , let \(C_k\) be the set of \(y \in C\) such that

\[
f _ {1} (y) <   f (y), f _ {2} (y) <   f (y), \dots , f _ {k - 1} (y) <   f (y), f _ {k} (y) = f (y).
\]

The \(C_{k}\) are disjoint convex cones with union C. Let \(x_{1}, x_{2}, \ldots, x_{n}\) in C be such that \(x_{1} + x_{2} + \cdots + x_{n} = x\) . Let \(y_{k}\) be the sum of the \(x_{i}\) that belong to \(C_{k}\) . Then \(y_{1} + y_{2} + \cdots + y_{p} = x\) . Since f is affine on \(C_{k}\) , \(f(y_{1}) + \cdots + f(y_{p}) = f(x_{1}) + \cdots + f(x_{n})\) . Therefore

\[
f (x) = \sup \bigl (f (y _ {1}) + \dots + f (y _ {p}) \bigr),\tag{19}
\]

where \((y_{1}, y_{2}, \ldots, y_{p})\) runs over the set of sequences of p points of C such that \(y_{1} + y_{2} + \cdots + y_{p} = x\) . Set \(D = C \cap (x - C)\) . Since D is compact (TVS, II, §6, No.~8, Cor. 2 of Prop. 11), so is the set of elements \((y_{1}, \ldots, y_{p})\) of \(D^{p}\) such that \(y_{1} + \cdots + y_{p} = x\) , thus the supremum (19) is attained.

Lemma 6. --- We maintain the hypotheses and notations of Lemma 5, and assume the \(f_{\lambda}\) to be positive. The function \(\overline{f}\) is positively homogeneous, concave and upper semi-continuous in C. It is affine if C is lattice-ordered.

It is clear that \(\overline{f}\) is positively homogeneous. Let x, y belong to C. If \(x_{1}, \ldots, x_{m}, y_{1}, \ldots, y_{n}\) in C are such that \(x_{1} + \cdots + x_{m} = x\) , \(y_{1} + \cdots + y_{n} = y\) , then \(x_{1} + \cdots + x_{m} + y_{1} + \cdots + y_{n} = x + y\) , therefore

\[
f (x _ {1}) + \dots + f (x _ {m}) + f (y _ {1}) + \dots + f (y _ {n}) \leqslant \overline {{f}} (x + y);
\]

it follows that \(\overline{f}(x) + \overline{f}(y) \leqslant \overline{f}(x + y)\) , therefore \(\overline{f}\) is concave. Let \(L\) (resp. \(L_{\lambda}\) ) be the set of \((t, x) \in \mathbf{R} \times E\) such that \(x \in C\) and \(0 \leqslant t \leqslant \overline{f}(x)\) (resp. \(0 \leqslant t \leqslant f_{\lambda}(x)\) ). Each of the \(L_{\lambda}\) is closed in the weakly complete proper convex cone \(\mathbf{R}_{+} \times C\) , therefore the sum \(\sum_{\lambda \in \Lambda} L_{\lambda}\) is closed (TVS, II, §6, No.~8, Cor. 2 of Prop. 11). By Lemma 5, this sum is equal to L. Therefore L is closed, which proves that \(\overline{f}\) is upper semicontinuous. Finally, assume C is lattice-ordered, and let us prove that \(\overline{f}\) is convex. Let \(x, y\) belong to C and let \(\varepsilon > 0\) . There exist \(z_1, z_2, \ldots, z_n\) in C such that \(f(z_1) + \cdots + f(z_n) \geqslant \overline{f}(x + y) - \varepsilon\) and \(z_1 + \cdots + z_n = x + y\) . The vector space C - C is lattice-ordered for the order induced by that of E (A, VI, §1, No.~9, Prop. 8). By the decomposition theorem (loc. cit., No.~10, Th. 1), there exist \(x_1, \ldots, x_n, y_1, \ldots, y_n\) in C such that

\[
x _ {1} + y _ {1} = z _ {1}, \dots , x _ {n} + y _ {n} = z _ {n}, x _ {1} + \dots + x _ {n} = x, y _ {1} + \dots + y _ {n} = y.
\]

Then, since \(f\) is positively homogeneous and convex,

\[
\begin{array}{r l} \overline {{f}} (x + y) & \leqslant \varepsilon + f (z _ {1}) + \dots + f (z _ {n}) \\ & \leqslant \varepsilon + f (x _ {1}) + f (y _ {1}) + \dots + f (x _ {n}) + f (y _ {n}) \\ & \leqslant \varepsilon + \overline {{f}} (x) + \overline {{f}} (y). \end{array}
\]

Since \(\varepsilon\) is an arbitrary number \(>0\) , we have proved that \(\overline{f}\) is indeed convex.

THEOREM 3 (Choquet). --- Let E be a Hausdorff weak locally convex space, C a weakly complete proper convex cone with vertex 0 in E, G the union of the extremal generators of C, K a compact convex subset of C, \(\lambda\) and \(\lambda'\) positive measures of mass 1 on K, admitting the same barycenter, such that \(\lambda^{*}(\mathrm{K} - (\mathrm{K} \cap \mathrm{G})) = \lambda'^{*}(\mathrm{K} - (\mathrm{K} \cap \mathrm{G})) = 0\) . Assume that C is lattice-ordered. Then, for every lower semi-continuous, positively homogeneous convex function \(f \geqslant 0\) on C, \(\lambda^{*}(f|K) = \lambda'^{*}(f|K)\) .

Let \(\mathcal{A}\) (resp. \(\mathcal{A}'\) ) be the set of restrictions to \(\dot{\mathrm{C}}\) of the continuous linear forms (resp. affine functions) on E. We know (TVS, II, §5, No.~4, Remark 2) that \(f\) is the upper envelope of the set of elements of \(\mathcal{A}\) that are \(\leqslant f\) . The set of functions of the form \(\sup(f_1, \ldots, f_p)\) , where \(f_1, \ldots, f_p\) belong to \(\mathscr{A}\) , \(f_1 \geqslant 0, \ldots, f_p \geqslant 0\) , is an increasing directed set and has \(f\) as its upper envelope. Taking into account §1, No.~1, Th. 1, it suffices to verify the equality \(\lambda(f|K) = \lambda'(f|K)\) when \(f\) is of the preceding form.

Define \(\overline{f}\) as in Lemma 5. It is clear that \(\overline{f}(y) = f(y)\) if \(y \in G\) . Since \(\lambda^{*}(\mathrm{K} - (\mathrm{K} \cap \mathrm{G})) = 0\) , we have \(\lambda(f|K) = \lambda(\overline{f}|K)\) . By Lemma 6, \(\overline{f}\) is affine and upper semi-continuous. Therefore \(\overline{f}|K\) is the lower envelope of a decreasing directed set of restrictions of elements of \(\mathcal{A}'\) to K (TVS, II, §5, No.~4, Prop. 6). Let \(x \in K\) be the barycenter of \(\lambda\) . If \(g \in \mathcal{A}\) then \(\lambda(g|K) = g(x)\) . Therefore \(\lambda(\overline{f}|K) = \overline{f}(x)\) (§4, No.~4, Cor. 2 of Prop. 5). Thus \(\lambda(f|K) = \overline{f}(x)\) , and one sees similarly that \(\lambda'(f|K) = \overline{f}(x)\) .

COROLLARY. --- Let E be a Hausdorff locally convex space, C a proper convex cone with vertex 0 in E, admitting a compact sole M, and let G be the union of the extremal generators of C. Let \(x \in M\) . If C is lattice-ordered, then there exists at most one positive measure \(\lambda\) of mass 1 on M, such that \(\lambda^{*}(\mathrm{M} - (\mathrm{G} \cap \mathrm{M})) = 0\) , and admitting x as barycenter.

Replacing the topology of E by the weakened topology (which does not change the topology of M), one can suppose E to be a weak space. Let \(\lambda\) and \(\lambda'\) be two measures on M having the stated properties, and let h be a continuous linear form on E such that M is the intersection of C and the hyperplane with equation \(h(x)=1\) . Let S be the subset of \(\mathcal{C}(\mathrm{M})\) consisting of the restrictions to M of the positively homogeneous and continuous convex functions \(\geqslant0\) on C. The cone C is weakly complete (TVS, II, §7, No.~3). By Th. 3, \(\lambda(f)=\lambda'(f)\) for every \(f\in S\) .

If \(f_{1}, f_{2}, f_{3}, f_{4}\) belong to \(\mathcal{S}\) , then

\[
\sup \left(f _ {1} - f _ {2}, f _ {3} - f _ {4}\right) = \sup \left(f _ {1} + f _ {4}, f _ {3} + f _ {2}\right) - \left(f _ {2} + f _ {4}\right) \in \mathscr {S} - \mathscr {S}
\]

\[
\inf \left(f _ {1} - f _ {2}, f _ {3} - f _ {4}\right) = - \sup \left(f _ {2} - f _ {1}, f _ {4} - f _ {3}\right) \in \mathscr {S} - \mathscr {S}.
\]

Since \(h|M \in \mathcal{S}\) , \(\mathcal{S} - \mathcal{S}\) contains the constant functions. If \(x\) and \(y\) are two distinct points of \(M\) , there exists a continuous linear form on \(E\) that separates \(x\) and \(y\) , and this form is the difference of two continuous linear forms that are positive on \(C\) (TVS, II, §6, No.~8, Lemma 1). It follows from the foregoing that for \(\alpha, \beta\) real, there exists \(f \in \mathcal{S} - \mathcal{S}\) such that \(f(x) = \alpha\) , \(f(y) = \beta\) . Then \(\mathcal{S} - \mathcal{S}\) is dense in \(\mathcal{C}(M)\) for the topology of uniform convergence (GT, X, §4, No.~1, Cor. of Prop. 2). Since \(\lambda\) and \(\lambda'\) coincide on \(\mathcal{S} - \mathcal{S}\) , we have \(\lambda = \lambda'\) .

